
\section{实验}
\label{appendix-exps}
\subsection{数据集}
在具有挑战性的nuScenes \cite{caesar2020nuscenes} 数据集上进行实验，该数据集包含训练集/验证集/测试集分别为700/150/150个场景。每个场景时长约20秒，包含约40个关键帧（以2Hz标注），每个样本包括6台相机拍摄的6张图像覆盖水平360°视场角，以及激光雷达和雷达传感器采集的点云。
\subsection{配置细节}

从结构上看，SparseAD由三部分组成：传感器编码器、稀疏感知和运动规划器。具体而言，\model 有两个变体：\model-Base和\model-Large分别采用VoVNetv2\cite{lee2020centermask} 和ViT-Large\cite{dosovitskiy2020image} 作为图像骨干网络，输入图像尺寸分别为 $1600\times640$ 和 $1600\times800$。

对于nuScenes数据集，SparseAD的范围在X和Y轴上均设置为 $\left[-51.2m, 51.2m\right]$。SparseAD中的所有子任务（包括检测、跟踪、在线建图、运动预测和规划）均在此范围内进行和评估。在SparseAD中，将EMMB的最大历史长度 $M$ 设置为5，对应约2.5秒的历史时序信息。

\subsubsection{稀疏感知}

对于检测\&跟踪任务，每帧设置644个可学习的检测查询，以及来自EMMB的256个场景级检测记忆查询（遵循StreamPETR\cite{wang2023exploring}），以及来自上一帧传播的高分查询作为跟踪查询。对于每个跟踪查询，设置0.5的丢弃概率，并在训练中以0.2的概率引入近似负样本以平衡实例级跟踪查询和场景级检测查询的监督。对于在线建图任务，分别使用 $\left[35, 30, 15, 25\right]$ 个可学习在线建图查询和 $\left[14, 12, 6,10\right]$ 个场景级建图记忆查询，对应\textit{分隔线}、\textit{人行横道}、\textit{道路段}和\textit{车道}，参照BeMapNet\cite{qiao2023end}。分段贝塞尔曲线 $\langle \text{最大分段数}, \text{最大阶数}\rangle$ 对
\textit{分隔线}、\textit{人行横道}、\textit{道路段}和\textit{车道}分别设置为
$ \langle3, 2\rangle, \langle1, 1\rangle, \langle7, 3\rangle$ 和 $ \langle7, 3\rangle$。

\subsubsection{运动规划器}

该模块利用最近5帧的时序信息（设置 $M=5$）预测其他智能体未来12帧（6秒）的运动，同时规划自车未来6帧（3秒）的运动。在nuScenes \texttt{train} 集上执行K-Means聚类算法，为运动查询获取先验意图点和轨迹。对于每个智能体（包括自车），设置6个运动查询以预测多模态未来轨迹。在训练过程中，仅监督与真实轨迹最接近的预测轨迹。自车的多模态运动预测作为强先验信息输入到规划优化器。最终，SparseAD为自车输出高质量、安全且可控的规划轨迹。

\subsubsection{训练与推理}
\label{sec:train_infer_detail}

SparseAD可以端到端地在线训练和测试。训练共分为三个阶段。在\textbf{第一阶段}，稀疏感知模块从零开始训练，获取驾驶场景的基本感知能力，包括检测和在线建图。在\textbf{第二阶段}，引入两层级记忆以优化稀疏感知模块的跟踪性能。在\textbf{第三阶段}，冻结骨干网络和稀疏感知模块，训练运动规划器模块。特别地，\textbf{第一阶段和第二阶段可以合并}进行训练，以缩短训练时间，但性能会有轻微下降，这是一个权衡。在实验中，使用AdamW\cite{loshchilov2017decoupled} 优化器和余弦退火\cite{loshchilov2016sgdr} 调度器在8张NVIDIA A100 80GB GPU上训练模型。特别是，如 \cref{tab:joint-results} 所示，SparseAD-Base在整个训练过程中仅需最多 \textbf{17.5GB} GPU内存，这意味着可以在\textbf{NVIDIA RTX 3090 24GB GPU}上完成整个训练过程。\cref{sec:experiments} 中的消融实验均在NVIDIA RTX 3090 24GB GPU上进行，输入图像尺寸设置为 $800\times320$。在推理过程中，当某智能体连续8帧（4秒）未被模型检测到时，对应的实例级记忆将被移除，以保持性能与效率之间的平衡。

\subsection{评估指标}
\subsubsection{稀疏感知指标}

为与其他流行方法保持一致，采用nuScenes \cite{caesar2020nuscenes} 数据集上感知任务评估常用的指标。具体而言，使用\textbf{NDS}（nuScenes检测分数）和\textbf{mAP}（平均精度）作为评估\model 检测性能的主要指标。作为nuScenes数据集自身提供的检测性能综合指标，NDS综合考量了不同维度的指标，包括\textbf{mAP}、\textbf{mATE}（平均平移误差）、\textbf{mAOE}（平均朝向误差）、\textbf{mASE}（平均尺度误差）、\textbf{mAVE}（平均速度误差）和\textbf{mAAE}（平均属性误差）。对于跟踪任务，采用\textbf{AMOTA}（平均多目标跟踪精度）、\textbf{AMOTP}（平均多目标跟踪精确度）、\textbf{RECALL}（召回率）和\textbf{IDS}（ID切换）等指标进行评估，其中AMOTA和AMOTP分别是不同召回率下\textbf{MOTA}和\textbf{MOTP}的平均结果。具体计算方法可参照 \cite{caesar2020nuscenes}。对于在线建图，将地图元素分为四种类型：分隔线、人行横道、道路段和车道，计算各类型在TP阈值为 $[0.5m, 1m, 2m]$ 下的\textbf{AP}（平均精度），最终获得每个类别的\textbf{mAP}（平均精度均值）作为评估指标。

\subsubsection{运动规划器指标}
与主流的运动预测方法保持一致，采用\textbf{minADE}（最小平均距离误差）和\textbf{minFDE}（最小最终距离误差）等指标作为评估SparseAD预测性能的主要指标。此外，还使用\textbf{MR}（缺失率）和\textbf{EPA}（端到端预测精度，End-to-end Prediction Accuracy）\cite{gu2023vip3d} 来评估端到端运动预测的质量，这有助于更全面地反映性能，因为minADE和minFDE指标仅在TP上计算。基准测试与VIP3D \cite{gu2023vip3d} 保持一致以确保公平比较。对于规划，利用先前方法 \cite{hu2023planning, jiang2023vad} 广泛使用的\textbf{L2}误差和\textbf{ColRate}（碰撞率）来展示自车的规划性能。具体而言，L2误差表示自车规划轨迹与实际轨迹之间的拟合程度，ColRate反映自车与周围智能体之间发生碰撞的概率。需要指出的是，重新审视了其他方法中碰撞率的计算方法，并采用稀疏方式更精确地计算碰撞率，详见 \cref{appendix-discussion}。

\subsection{补充实验}
\begin{table}[t]
\centering
\caption{nuScenes \texttt{test} 数据集上的3D检测和多目标跟踪结果。}
\label{tab:det_tracking_test_set}
\tiny
\resizebox{\textwidth}{!}{
% {
\setlength{\tabcolsep}{3.5pt}
\begin{tabular}{l|c|c|ccc|l|cccc}
\toprule
\textbf{方法}&  \textbf{mAP}$\uparrow$  &\textbf{NDS}$\uparrow$  & \textbf{mATE}$\downarrow$  &\textbf{mAOE}$\downarrow$   &\textbf{mAVE}$\downarrow$ & \textbf{方法}& \textbf{AMOTA}$\uparrow$  & \textbf{AMOTP}$\downarrow$ & \textbf{RECALL}$\uparrow$ & \textbf{IDS}$\downarrow$\\
\midrule
\multicolumn{5}{l}{\textit{纯检测方法}}&&\multicolumn{5}{l}{\textit{纯跟踪方法}}\\
\midrule
% ResNet101
StreamPETR\cite{wang2023exploring}   &0.546& 0.633 &0.485&  \underline{0.307}& 0.250                  &DEFT\cite{chaabane2021deft} &0.177 &1.564 &0.338&  6901\\
PETRv2\cite{liu2023petrv2}         &0.490& 0.582& 0.561& 0.361 &0.343               &QD3DT\cite{hu2022monocular} &0.217& 1.550 &0.375  &6856\\ % PETR Table 1
BEVFormer \cite{li2022bevformer}    & 0.481& 0.569& 0.582&  0.375& 0.378            & CC-3DT$^*$~\cite{fischer2022cc} &0.410 &1.274 &0.538&3334\\ % BEVFormer Table 4.
Sparse4D\cite{lin2022sparse4d}  &0.511 & 0.595&0.533&  0.369 &0.317                 &QTrack\cite{yang2022quality}& 0.480 &1.107& 0.569&  1484\\ % PolarDETR Table 1
SRCN3D\cite{shi2022srcn3d}&0.396 &0.463&0.673 & 0.403& 0.875& SRCN3D\cite{shi2022srcn3d}& 0.398 &1.317 &0.538&-\\
DORT\cite{qing2023dort}& 0.545& 0.625&\underline{0.440} & 0.370& 0.250 &DORT\cite{qing2023dort}&\underline{ 0.576} &\underline{0.951}&-&-\\
\midrule
\multicolumn{7}{l}{\textit{端到端多目标跟踪方法}}\\
\midrule
MUTR3D\cite{zhang2022mutr3d}  &-&-&-&-&-& MUTR3D\cite{zhang2022mutr3d} &0.270& 1.494 &0.411& 6018\\
PF-Track\cite{pang2023standing} &-&-&-&-&- &PF-Track\cite{pang2023standing} &0.434 &1.252 & 0.538&  \textbf{249}\\
DQTrack\cite{li2023end}&-&-&-&-& -& DQTrack\cite{li2023end} &0.523& 1.096 &  0.622 &1204\\
Sparse4Dv3\cite{lin2023sparse4d} &\textbf{0.570}& \textbf{0.656}&\textbf{0.412} & 0.312& \textbf{0.210} & Sparse4Dv3\cite{lin2023sparse4d} &0.574& 0.970  &\underline{0.669}  &669\\
\midrule
\multicolumn{7}{l}{\textit{端到端自动驾驶方法}}\\
\midrule
% UniAD\cite{hu2023planning} &-&-&-	&	-	&-& UniAD\cite{hu2023planning} &- &- &- &-\\
\rowcolor[gray]{.9}  \model-B& 0.514 & 0.605 &0.505 &0.337 &0.308	& \model-B &0.502 &1.065 &	0.643 &349 \\
\rowcolor[gray]{.9}  \model-L& \underline{0.566} & \underline{0.648} &0.454 &\textbf{0.288} &\underline{0.237}	& \model-L &\textbf{0.588} &\textbf{0.904} &	\textbf{0.691} &\underline{344} \\
\bottomrule
\end{tabular}}
\end{table}

\subsubsection{测试集评估}
尽管 \cref{sec:experiments} 中的实验已在nuScenes \texttt{val} 集上详细评估了SparseAD的综合性能并与其他方法进行了比较，本文也报告了SparseAD在nuScenes \texttt{test} 集上的性能以进一步验证感知性能。如 \cref{tab:det_tracking_test_set} 所示，\model-Base与流行方法 \cite{liu2023petrv2, li2022bevformer, lin2022sparse4d, shi2022srcn3d} 相比取得了可比性能，但仍逊于当前最优的单任务方法 \cite{lin2023sparse4d, wang2023exploring}。与此同时，\model-Large达到了 \textbf{56.6\%} mAP、\textbf{64.8\%} NDS和 \textbf{58.8\%} AMOTA，甚至略优于SOTA方法\cite{lin2023sparse4d}，展示了端到端范式在自动驾驶领域的巨大潜力。值得注意的是，UniAD\cite{hu2023planning} 未在nuScenes \texttt{test} 集上进行评估，因此没有可以完全公平比较的方法。作为端到端方法，仅可在nuScenes官方基准上评估检测和跟踪结果以与其他方法公平比较。因此，无法提供其他任务（如建图、运动预测和规划）在nuScenes \texttt{test} 集上的评估结果。

\begin{table*}[t]
		\caption{SparseAD-B在第三阶段训练中带或不带端到端梯度优化的综合性能对比。B.：骨干网络，P.：稀疏感知，M.：运动规划器，T.Mem：训练内存占用。}
		\begin{center}
			\tiny
			% \renewcommand\arraystretch{1.5}
			\vspace{-10pt}
			\resizebox{1\textwidth}{!}
			{
				\begin{tabular}{ccc|cc|cc|ccc|cccc|cccc|c}
					\toprule
					\multicolumn{3}{c|}{梯度} &
					\multicolumn{2}{c|}{检测} &
					\multicolumn{2}{c|}{跟踪} &
					\multicolumn{3}{c|}{运动预测} &
					\multicolumn{4}{c|}{规划L2$\downarrow$} &
					\multicolumn{4}{c|}{规划碰撞率$\downarrow$} &
					效率 \\
					B.&P.&M.& mAP$\uparrow$ & NDS$\uparrow$ & AMOTA$\uparrow$&Recall$\uparrow$& minADE$\downarrow$&minFDE$\downarrow$&MR$\downarrow$ &1s &2s &3s &平均&1s&2s&3s &平均&  训练内存 \\
					\midrule
					&&\Checkmark& 0.475& 0.578& 0.530& 0.608 & 0.83& 1.58& 0.187&0.15&0.31&0.57& 0.35&0.00&0.06&0.21& 0.09&4.0\\
					\Checkmark&\Checkmark&\Checkmark&0.455& 0.559& 0.521 & 0.622&0.83&	1.58&	0.180&0.15&0.31&0.56& 0.34 &0.00&0.04&0.15&0.08&18.1\\
					\midrule
					\multicolumn{3}{c|}{\textit{差值}} &\textcolor{blue}{-0.020} & \textcolor{blue}{-0.019}& \textcolor{blue}{-0.009} & \textcolor{red}{+0.014} & 0.00 & 0.00 & \textcolor{red}{-0.007}& 0.00&0.00&\textcolor{red}{-0.01}&\textcolor{red}{-0.01}&0.00&\textcolor{red}{-0.02}&\textcolor{red}{-0.06}&\textcolor{red}{-0.01}& \textcolor{blue}{+12.1} \\

					\bottomrule
				\end{tabular}% }
			}
		\end{center}
		% \vspace{2pt}
		% \vspace{-20pt}
		\label{tab:joint-results-gradient}
	\end{table*}

\subsubsection{端到端梯度优化。} 需要注意的是，在\textbf{第三阶段}训练过程中，SparseAD并未开启整个模型的梯度。然而，由于SparseAD的高效设计，实现整个模型的端到端梯度优化是可行的。基于SparseAD-Base，评估了第三阶段训练中不同梯度设置对性能的影响。如 \cref{tab:joint-results-gradient} 所示，运动规划器的梯度传播到上游感知任务时，运动预测和规划的子任务性能略微提升，端到端优化带来的额外内存占用对学术级计算能力（如NVIDIA RTX 3090 GPU）也是可接受的。然而，端到端优化也略微降低了SparseAD的感知能力。考虑到自车规划的准确性和安全性是端到端自动驾驶方法的最终目标，感知性能的轻微下降在此处是可接受的。总体而言，本文认为第三阶段的端到端优化可以被视为感知与规划性能之间的一种权衡。

\subsection{运行时间分析}

\begin{table*}[t]
		\caption{模块运行时间。推理速度在单张NVIDIA RTX A100 80GB GPU上对SparseAD-Base和SparseAD-Large进行测量。}
		\begin{center}
			\tiny
			\vspace{-10pt}
			{
				\begin{tabular}{p{4cm}|p{2cm}<{\centering}p{2cm}<{\centering}|p{2cm}<{\centering}p{2cm}<{\centering}}
					\toprule
					\multirow{2}{*}{模块}&
					% \multicolumn{2}{c|}{VoVNet-V2-99\cite{lee2020centermask}} &
					% \multicolumn{2}{c}{ViT-Large\cite{dosovitskiy2020image}}  \\
	                \multicolumn{2}{c|}{SparseAD-B} &
					\multicolumn{2}{c}{SparseAD-L}  \\
					& 延迟(ms) & 占比 & 延迟(ms) & 占比\\
					\midrule
					Backbone& 85& 29.6\%&1153&83.1\%\\
					Sparse Perception&104& 36.2\%&130&9.4\%\\
					Motion Planner&75&26.1\%&75&5.4\%\\
					Output\&Memory Update&23 &8.1\%&28 & 2.1\%\\
					\midrule
					\textit{总计}&287 & 100.0\% & 1386 & 100\%\\
					\bottomrule
				\end{tabular}% }
			}
		\end{center}
		\label{tab:runtime}
	\end{table*}

评估了SparseAD-Base和SparseAD-Large各模块的运行时间，结果如 \cref{tab:runtime} 所示。SparseAD-Base的推理延迟为 \textbf{287ms}，其中稀疏感知是最耗时的模块，占总延迟的36.2%。值得注意的是，SparseAD中维护EMMB的运行时间为23ms，仅占总延迟的约8.1%，这使得能够高效利用长时间时序信息。另一方面，SparseAD-Large的推理延迟为 \textbf{1386ms}，其中庞大的ViT-Large骨干网络成为推理延迟的主要组成部分，占总推理延迟的83.1%。SparseAD-Large相比SparseAD-Base在综合性能上有所提升，而稀疏感知、运动规划器和输出\&内存更新方面的延迟仅有轻微增加。这些结果也表明SparseAD的设计能够高效地从先进的大规模骨干网络中获益。
